\section{Day 7: Ideals of Localization and Sheaf properties (Sep.\ 29, 2026)}
\begin{theorem} \label{thm:localization_ideals}
    We have the following properties of localization,
    \begin{enumerate}[(i)]
        \item $\kp R_\kp$ is the \textit{unique} maximal ideal (i.e., $R_\kp$ is a local ring).
        \item $\abs{\Spec R_\kp} \cong \{\kq R_\kp \mid \kq \in \Spec R, \, \, \kq \subset \kp\} \subset \Spec R$, i.e., a homeomorphism.
    \end{enumerate}
\end{theorem}
\begin{proof}
    $R_\kp$ considered as germs $(f/g^i, D(g))$ has an obvious ideal: the set of germs vanishing at $\kp$. This is an ideal, since, if we have a germ $(f/g^i, D(g))$ not in this ideal, then restricting further to $D(fg) = \kp$ yields us a unit
    \[ \left(\frac{f}{g^i}\biggr\rvert_{D(fg)}, D(fg)\right), \]
    so the germ is a unit. Consequently, no proper ideal can contain an element outside of this ideal, so this ideal is maximal.

    Finally, note that any element of this ideal $(f/g^i, D(g))$ has $f(\kp) = 0$, and hence is equivalent, up to a unit, to $(f, \Spec R)$ with $f \in \kp$. Thus, this ideal lies in $\kp R_\kp$. The converse is obvious, since, if $f \in \kp$, then $\kp R_\kp$ is in the ideal, and we have $f(\kp) = 0$.

    For (ii), we claim that each prime ideal is of the form $\kq R_\kp$ for some $\kq \in \Spec R$. As shown in (i), any ideal in $R_\kp$ is of the form $I R_\kp$ for some $I \subset R$. In particular, this ideal has the further property that if $rs \in I$ and $s \notin \kp$, then $r \in I$. Given an ideal with this property, we want to show that $I R_\kp$ is prime is equivalent to $I \subset R$ being prime.

    First, if $I R_\kp$ is prime, suppose $ab \in I$; we want to show that one of $a, b$ is in $I$. Indeed, note that, if $a_\kp b_\kp \in I R_\kp$ (where $a_\kp, b_\kp \in R_\kp$), then one of $a_\kp$ or $b_\kp$ lie in $I R_\kp$; suppose it is $a_\kp$. But then $a_\kp \in I R_\kp$ implies that there exists $f \in \kp$ such that $fa \in I$, so $a \in I$. For the converse, we have that
    \[ \frac{R_\kp}{I R_\kp} \cong \left(\frac{R}{I}\right)_\kp \to \Frac(R/I). \]
    
    Finally, given $\alpha = (f/g^i, D(g)) \in R_\kp$, we want to show that $D(\alpha)$ is the intersection of an open subset of $\Spec R$ with $\Spec R_\kp$. Indeed, $D(f) \cap \Spec R_\kp = D(\alpha)$.
\end{proof}

We now talk about the sheaf property. Consider $C^0(M, \RR)$, and consider the open cover $\{U_\alpha\}_{\alpha \in \SA}$. We may evaluate continuous functions on each $U_\alpha$, and we may look at continuous functions on each $U_\alpha$ that are compatible on each $U_\alpha \cap U_\beta$, i.e.,
\[ \restr{f_\alpha}{U_\alpha \cap U_\beta} = \restr{f_\beta}{U_\alpha \cap U_\beta}, \]
to patch together globally continuous functions on the entire of $M$; this is because continuity is a local property.

As an example, let us consider $\Spec \ZZ$. $D(9)$ and $D(5)$ cover $\Spec \ZZ$, since the ideal generated by $9$ and $5$ contain $1$. In terms of the spectrum,
\[ D(9) = \Spec \ZZ \setminus (3), \quad D(5) = \Spec \ZZ \setminus (5). \]
Clearly, the union is the whole of $\Spec \ZZ$, since each prime is either divisible $3$ or $5$, but not both (usually, it's divisible by neither). The intersection of $D(9)$ and $D(5)$ is $D(45) = \Spec \ZZ \setminus (3, 5)$. 

\newpage
Functions in $D(9)$ are of the form $\frac{x}{9^i}$ and in $D(5)$, they're of the form $\frac{y}{5^j}$. Assuming
\[ \frac{x}{9^i} = \frac{y}{5^j} \in \ZZ_{45} \subset \QQ, \]
we have that
\[ \frac{x}{9^i} = \frac{y}{5^j} = z \]
for some unique integer $z$. $D(9^i) = D(9)$ and $D(5^j) = D(5)$, so $(1) = (9^i, 5^j)$, so $(1) = (9, 5)^{i+j}$, and we have $1 = c9^i + d5^j$.\footnote{``we now do a really weird property of fractions that's only taught in russia''} Observe that
\[ \frac{x}{9^i} = \frac{y}{5^j} = \frac{cx + dy}{c 9^i + d 5^j} = cx + dy. \]
Let us look at a more general example. In $\RR[x, y]$, we have that
\[ Z(y - x^2) \leftrightarrow Z(y^2 - x^4 + 1) \subset Z(y - x^2) \cup Z(y^7 - x^4 + 1). \]
$Z(y - x^2)$ and $Z(y^2 - x^4 + 1)$ are disjoint; how do we glue them together? By taking $1 = (y^2 - x^4 + 1) - (y + x^2)(y - x^2)$... and then we ran out of time. Hunter will continue this example next lecture.